\color{unidarkgreen}{\subsection{Технологические защиты}}
\color{black}

\subsubsection{Общие сведения}

\begin{enumerate}[label=\arabic{section}.\arabic{subsection}.\arabic{subsubsection}.\arabic*,  align=left, labelsep=4pt, leftmargin=0pt, itemindent=57pt]

\item
Технологические защиты от повышения давления, температуры масла и обмоток предназначены
для предотвращения повреждения оборудования вследствие нарушения режима его работы,
обеспечивая надёжную и безопасную работу трансформатора. Для трансформатора
предусматриваются функциональные блоки приема и обработки сигналов от датчиков температуры
масла и обмотки, дополнительно для ЛРТ предусматривается прием и обработка сигналов от реле
давления (рисунок \ref{pic:TZT}). 

\item
Предусмотрена возможность приема сигналов срабатывания сигнального контакта ДТм Т (ДТо Т),
отключающего контакта ДТм Т (ДТо Т) и низкой изоляции цепи отключающего контакта ДТм Т (ДТо Т), по
разным цепям от двух источников ПДС1 и ПДС2. Для приема сигналов от ПДС1 предусмотрены входные
сигналы \textit{«Ср. сигн контакта ДТм Т (ДТо Т) 1 цепь»}, \textit{«Ср. авар контакта ДТм Т (ДТо Т) 1 цепь»} и
\textit{«Низкая изол. ДТм Т (ДТо Т) авар 1 цепь»}, а для приема сигналов от ПДС2 предусмотрены входные
сигналы \textit{«Ср. сигн контакта ДТм Т (ДТо Т) 2 цепь»}, \textit{«Ср. авар контакта ДТм Т (ДТо Т) 2 цепь»} и
\textit{«Низкая изол. ДТм Т (ДТо Т) авар 2 цепь»}. Если сигналы срабатывания ДТм Т (ДТо Т) заводятся
непосредственно со шкафа ШОМ силового трансформатора конфигурируются только входные сигналы
\textit{«Ср. сигн контакта ДТм Т (ДТо Т) 1 цепь»}, \textit{«Ср. авар контакта ДТм Т (ДТо Т) 1 цепь»} и \textit{«Низкая изол.
ДТм Т (ДТо Т) авар 1 цепь»}, а действие ДТм Т (ДТо Т) по второй цепи сигнального и аварийного контакта
выводиться из работы с помощью уставок \textbf{«ДТм Т (ДТо Т) сигн 2 цепь»} и \textbf{«ДТм Т (ДТо Т) авар 2 цепь»}.

\item
Ввод в работу ДТм Т (ДТо Т) осуществляется с помощью уставок \textbf{«ДТм Т (ДТо Т) сигн 1 цепь»},
\textbf{«ДТм Т (ДТо Т) сигн 2 цепь»}, \textbf{«ДТм Т  (ДТо Т) авар 1 цепь»} и \textbf{«ДТм Т (ДТо Т) авар 2 цепь»}.
Предусмотрена возможность действия срабатывания аварийного контакта ДТм (ДТо) на отключение
трансформатора с помощью уставки \textbf{«Откл от ДТм Т (ДТо Т) авар»}.

\item
Оперативный ввод/вывод ТЗ производится с помощью БУ (блока управления функциями) \textit{«Вывод ТЗ Т»}.

\item
Непрерывный контроль изоляции цепи аварийного контакта ДТм (ДТо) осуществляется с помощью
реле контроля изоляции (РКИ). В случае снижения сопротивления изоляции РКИ подает сигнал \textit{«Низкая
изол. ДТм Т (ДТо Т) авар 1 цепь»} (\textit{«Низкая изол. ДТм Т (ДТо Т) авар 2 цепь»}) для блокировки
срабатывания ДТм Т (ДТо Т) от аварийного контакта. Время срабатывания блокировки от сигналов РКИ
задается с помощью уставки \textbf{«Тбл ТЗ Т»}. Ввод блокировки ДТм Т (ДТо Т) от РКИ осуществляется с
помощью уставки \textbf{«Блок ТЗ Т»}. Для сброса блокировки от РКИ, используется команда управления
\textit{«Возврат ТЗ Т»}.

\item
Технологическая защита ДТм Т аналогична ДТо Т, ДТм ЛРТ, ДТо ЛРТ. На рисунке \ref{pic:TZT}
представлен Функциональный блок ТЗ Т.

\begin{figure}[H]
\centering
\includegraphics[width=0.82\textwidth,height=0.82\textheight,keepaspectratio]{\fbpath/011.1061 ТЗ Т/img/TZT.pdf}
\captionof{figure}{Функциональный блок ТЗ Т}\label{pic:TZT}
\end{figure}


\small
\begin{longtable}{|>{\centering\arraybackslash}m{5.3cm}|>{\centering\arraybackslash}m{4.2cm}|>{\centering\arraybackslash}m{1cm}|>{\centering\arraybackslash}m{3.3cm}|>{\centering\arraybackslash}m{1.8cm}|}
% Первая страница: полный заголовок
\caption{Перечень уставочных параметров ТЗ\hfill\vspace{-0.5\baselineskip}}\label{desc:tbl1}\\
\hline
\rowcolor{gray!30}
Наименование уставки & Диапазон & Шаг & Значение по умолчанию & Единица измерения \\
\hline
\endfirsthead

% Промежуточные страницы: «Продолжение таблицы»
\caption*{\hspace{3pt}\emph{Продолжение таблицы \ref{desc:tbl1}}\hfill\ }\vspace{-0.5\baselineskip} \\
\hline
\rowcolor{gray!30}
Наименование уставки & Диапазон & Шаг & Значение по умолчанию & Единица измерения \\
\hline
\endhead

%===m>ДТЗ|-
%===m
\end{longtable}
\normalsize	

\end{enumerate}